\documentclass[10pt]{article}
\usepackage[T1]{fontenc}
\usepackage[utf8]{inputenc}
\usepackage[french]{babel}
\usepackage{amsmath,amssymb,mathrsfs,mathpazo}
\usepackage{geometry}
\geometry{a4paper,margin=2cm}
\usepackage{graphicx,array,tabularx,tabularray,colortbl,multirow,multicol,booktabs,enumitem,pifont,ulem,fancybox,diagbox,cancel}
\usepackage{pstricks,pstricks-add,pst-plot,pst-tree,pst-node,pst-text,pst-func,pst-math,pst-bspline,pst-3dplot,pst-eucl,pst-solides3d}
\let\pscalculate\undefined
\usepackage{pst-bezier}
\usepackage{tikz,pgfplots,tkz-tab}
\usetikzlibrary{arrows,arrows.meta,patterns,calc,positioning,trees,decorations.markings,shapes.geometric,shapes.misc,shapes.arrows,intersections,decorations.pathmorphing,decorations.pathreplacing,backgrounds,shadows,babel}
\pgfplotsset{compat=1.18}
\usepgfplotslibrary{fillbetween}
\usepackage[np]{numprint}
\usepackage{listings,makecell,pas-tableur,twoopt}
\let\barre\undefined
\usepackage[pstricks]{tablvar}
\definecolor{darkgreen}{RGB}{0,100,0}
\providecommand{\red}{\color{red}}
\providecommand{\blue}{\color{blue}}
\providecommand{\green}{\color{green}}
\providecommand{\black}{\color{black}}
\providecommand{\white}{\color{white}}
\providecommand{\gray}{\color{gray}}
\providecommand{\cyan}{\color{cyan}}
\providecommand{\magenta}{\color{magenta}}
\providecommand{\yellow}{\color{yellow}}
\providecommand{\violet}{\color{violet}}
\providecommand{\orange}{\color{orange}}
\providecommand{\euro}{\texteuro}
\providecommand{\EUR}{\texteuro}
\providecommand{\ieme}{\textsuperscript{e}}
\providecommand{\degres}{\ensuremath{^\circ}}
\providecommand{\degree}{\ensuremath{^\circ}}
\providecommand{\vv}[1]{\overrightarrow{#1}}
\providecommand{\Coord}[1]{#1}
\providecommand{\decosix}{}
\providecommand{\starredbullet}{\ensuremath{\bullet}}
\pagestyle{empty}
\setlength{\parindent}{0pt}
\usepackage{adjustbox}\DeclareUnicodeCharacter{2236}{:}\DeclareUnicodeCharacter{2190}{\ensuremath{\leftarrow}}
\providecommand{\R}{\mathbb{R}}
\providecommand{\N}{\mathbb{N}}
\providecommand{\Z}{\mathbb{Z}}
\providecommand{\Q}{\mathbb{Q}}
\providecommand{\D}{\mathbb{D}}
\providecommand{\C}{\mathbb{C}}
\providecommand{\e}{{\mathrm{e}}}
\providecommand{\vect}[1]{\overrightarrow{#1}}
\providecommand{\vectt}[1]{\overrightarrow{\mathrm{#1}}}
\providecommand{\vvt}[1]{\overrightarrow{\mathrm{#1}}}
\providecommand{\np}[2][]{{#2~#1}}
\providecommand{\vv}[1]{{\overrightarrow{#1}}}
\renewcommand{\d}{{\mathrm{d}}}
\providecommand{\degres}{{^\circ}}
\providecommand{\degre}{{^\circ}}
\providecommand{\blue}{\color{blue}}
\providecommand{\red}{\color{red}}
\providecommand{\enskip}{\quad}
\providecommand{\llbracket}{\lbrack\!\lbrack}
\providecommand{\rrbracket}{\rbrack\!\rbrack}
\providecommand{\euro}{{\text{€}}}
\providecommand{\no}{{\text{n°}}}
\providecommand{\interval}[2]{{[#1;#2]}}
\providecommand{\intervFO}[2]{{[#1;#2[}}
\providecommand{\dotfill}{\ldots}

\providecommand{\ds}{\displaystyle}
\providecommand{\textding}[1]{#1}
\providecommand{\Oij}{\(\mathrm{O};\vec{\imath},\vec{\jmath}\)}
\providecommand{\Oijk}{\(\mathrm{O};\vec{\imath},\vec{\jmath},\vec{k}\)}
\providecommand{\Ouv}{\(\mathrm{O};\vec{u},\vec{v}\)}
\providecommand{\intervFF}[2]{\ensuremath{[#1;#2]}}
\renewcommand{\euro}{\ifmmode\text{\texteuro}\else\texteuro\fi}
\DeclareUnicodeCharacter{00B0}{\ifmmode^{\circ}\else\textdegree\fi}
\DeclareUnicodeCharacter{20AC}{\ifmmode\text{\texteuro}\else\texteuro\fi}
% Énoncés : Éducation nationale. Transcription et corrigés : APMEP, Denis Vergès et contributeurs. Découpage et classement : CoopMaths, Sébastien Lozano, Éric Elter et contributeurs. Adaptation de la présentation : Maths Explorer.
% CC BY-SA : https://coopmaths.fr/www/about/
% https://coopmaths.fr/alea/static/e3c/2020/tex/e3c_2020_00_sujet12_2_cor.tex
\begin{document}
\noindent\textit{APMEP / CoopMaths. CC BY-SA. Présentation adaptée par Maths Explorer.}\par\medskip
\subsection*{1. a. Calculer $u_1$ et $u_2$.}
	
	\[
	u_1 = 300 \times 1,05 + 15 = 315 + 15 = 330
	\]
	\[
	u_2 = 330 \times 1,05 + 15 = 346,5 + 15 = 361,5
	\]
	
	\subsection*{b. Montrer que la suite $(u_n)$ ainsi définie, n'est ni arithmétique ni géométrique.}
	
	\[
	u_1 - u_0 = 30 \quad \text{et} \quad u_2 - u_1 = 31,5 \quad \Rightarrow \quad \text{la suite n'est pas arithmétique.}
	\]
	\[
	\dfrac{u_1}{u_0} = \dfrac{330}{300} = 1,1 \quad \text{et} \quad \dfrac{u_2}{u_1} = \dfrac{361,5}{330} \approx 1,095 \quad \Rightarrow \quad \text{la suite n'est pas géométrique.}
	\]
	

	
	\subsection*{2. On considère la suite $(v_n)$, définie pour tout entier naturel $n$, par : $v_n = u_n + 300$.}
	
	\subsubsection*{a. Calculer $v_0$, puis montrer que la suite $(v_n)$ est géométrique de raison $q = 1,05$.}
	
	\[
	v_0 = u_0 + 300 = 300 + 300 = 600
	\]
	\[
	v_{n+1} = u_{n+1} + 300 = 1,05u_n + 15 + 300 = 1,05u_n + 315 = 1,05(u_n + 300) = 1,05v_n
	\]
	
	\subsubsection*{b. Pour tout entier naturel $n$, exprimer $v_n$ en fonction de $n$, puis montrer que $u_n = 600 \times 1,05^n - 300$.}
	
	L'égalité $v_{n+1} = 1,05v_n$ vraie pour tout naturel $n$, montre que la suite $(v_n)$ est géométrique de raison $q = 1,05$ de premier terme 600.
	
	On sait qu'alors, quel que soit $n \in \mathbb{N}$, $v_n = 600 \times 1,05^n$.
	
	\subsection*{3. Est-il correct d'affirmer que la surface envahie par les chardons aura doublé au bout de 8 semaines ? Justifier la réponse.}
	
	On a $v_8 = 600 \times 1,05^8 \approx 886,473$.
	
	Or $v_8 = u_8 + 300 \iff u_8 = v_8 - 300 \approx 886,473 - 300 \approx 536,5$ soit moins du double de la surface initiale.

\end{document}
